% =====================================================================================
% Section 7 -- Discussion and limitations (restored for arXiv: docs/intent/paper-focus.md 18)
% No new numbers here.
% =====================================================================================
\section{Discussion and Limitations}
\label{sec:discussion}

\paragraph{Why decision fine-tuning does not make a reranker.}
A reranker fine-tune is trained on the very task it is tested on: telling relevant from non-relevant passages of a
query, with a ranking loss over hard negatives. The decision fine-tunes we tested learned to answer many kinds of
questions about many kinds of states, and Kev's training data contain no MS~MARCO and no BEIR~\cite{KevCards}. That
general skill transfers to relevance about as well as the untuned model's own instruction following, which is what we
observe up to 9B. The zero-shot subset shows that the gap is not about having seen the test domain: it persists on
collections that neither kind of fine-tune was trained on.

\paragraph{Why 27B is different.}
Kev-27B is the only decision fine-tune that adds to its untuned model. It differs from the smaller Kevs in more than
size: it is a full fine-tune of the released chat model, whereas the smaller Kevs are LoRA fine-tunes of pretrained
checkpoints, and it was trained on longer states (Sect.~\ref{sec:reranking}). We cannot tell which of these matters.
Size alone is not enough: Clef, a decision fine-tune of the same 27B model, adds nothing. Reranker fine-tuning also
adds less at 27B than on the smaller backbones, because the untuned 27B model already ranks well.

\paragraph{What calibration means for relevance.}
A probability of relevance is calibrated only with respect to a definition of relevance and its base rate. Our
results show three ways to move it: the wording of the question (Sect.~\ref{sec:calibration}), the labels of the
training data (our reranker fine-tunes) and the labels a scaling is fitted on (Platt scaling on MS~MARCO dev,
Sect.~\ref{sec:calibration}). A model that is ``calibrated out of the box'' is calibrated against a notion of
relevance that someone chose, and the user of a decision model has to put that notion into words. For ranking none of
this matters; for a threshold or a count it matters a lot.

\paragraph{Limitations.}
We compare released models, and the fine-tunes of a backbone differ in their training data as well as in their
objective. Our results therefore describe decision fine-tuning as it is available today, not RLCD as such, which is
unpublished and may be done better. The controlled comparison rests on our own reranker fine-tunes at 9B and 27B,
which were trained once, with one recipe and without any hyperparameter search. We measured calibration only on TREC
DL, where non-relevant passages are judged, and at a single relevance grade and a single threshold. Our question was
fixed before the runs but not chosen on development data, and we tried only one other wording. All decision models
are at most a month old, and newer versions may change the picture; we release all raw outputs so that the comparison
can be repeated.
